\chapter{PENUTUP}
\label{chap:penutup}

% Ubah bagian-bagian berikut dengan isi dari penutup

\section{Kesimpulan}
\label{sec:kesimpulan}

% Berdasarkan hasil penelitian yang telah dilakukan, maka dapat diambil beberapa kesimpulan sebagai berikut:

% \begin{enumerate}[nolistsep]

%   \item Dengan melakukan pendekatan non-linear menggunakan \emph{Machine Learning Multi Layer Perceptron} untuk kalibrasi kamera, maka hasil kalibrasi kamera lebih baik dibandingkan dengan menggunakan regresi polynomial sederhana.

%   \item Sistem kalibrasi dapat meng-kalibrasi semua arah pada kamera \emph{omnivision} tanpa perlu instalasi ulang hardware jika ada masalah.

% \end{enumerate}

Berdasarkan hasil penelitian yang telah dilakukan, didapat kesimuplan bahwa Metode Kalibrasi menggunakan \emph{Machine Learning} lebih baik daripada metode Kalibrasi menggunakan Regresi Polinomial. Hal itu dapat dilihat dari sisi akurasi dan presisi yang dihasilkan oleh metode \emph{Machine Learning} lebih baik daripada metode regresi polinomial. Error akurasi metode \emph{Machine Learning} sebesar 10.84 cm sedangkan metode regresi polinomial sebesar 20.77 cm. Error presisi metode \emph{Machine Learning} sebesar 1.20 cm dan 4.10 cm sedangkan metode regresi polinomial sebesar 10.01 cm dan 11.32. Dengan menggunakan metode \emph{Machine Learning} pada kalibrasi kamera \emph{Omnivision}, maka robot dapat bergerak dengan lebih baik. 

\section{Saran}
\label{chap:saran}


Untuk pengembangan lebih lanjut pada penelitian ini, bisa dilakukan beberapa hal antara lain:

\begin{enumerate}[nolistsep]

  \item Mempercepat proses pengambilan data. Hal ini bisa dilakukan dengan cara mengambil data secara otomatis yaitu menggantikan peran manusia yang meng-klik gambar kamera lalu memasangkan penanda satu persatu menjadi sistem otomatis dengan cara mendeteksi warna penanda yang ada. 

  \item Memperbaiki pola pengambilan data sehingga didapat data-data yang lebih baik untuk dilakukan proses \textit{Machine Learning}. Hal ini bisa dilakukan dengan cara menambah variasi data yang diambil, menambah jarak minimal dan maksimal data yang diambil, dan menyesuaikan persebaran data pada kamera sehingga hasil dari proses \emph{Machine Learning} lebih baik.
\end{enumerate}

